\subsection{The ordinary-category criterion and the domain of interchange}

\begin{result}{A.14}{Proved}{Deriving ordinariness from the two local inputs}
For every stable $(g,n)$, the asphericity and meridian injection in
A.5 imply that every exit mapping space is a disjoint union of
contractible components. This gives the application of EXIT\_ONE\_TYPE
P1b1 requested in added task 1, before invoking Jansen's global
identification. That identification supplies the explicit transporter
formula and its compatibility over all incidences.
\end{result}
\begin{proof}
For a fixed incidence of strata $S,T$, the space of exits with chosen
endpoints is the homotopy fibre of the relative link evaluation
$L_{S,T}\to S\times T$. The local model uses a positive radial
coordinate for each smoothed node; varying those positive radii is
contractible. Keeping their angular coordinates gives the torus
bundle used in A.5. Endpoint paths on $S$ and $T$ give the indicated
homotopy fibre; subdivision of these paths and the same normal
coordinates identify this construction on overlapping charts. This
is the mapping-space calculation of P1a1 in the supplied return,
with its two endpoint maps retained.

Both endpoints are aspherical by A.5. The link is aspherical too:
it is a torus bundle over the incidence cover of $S$, and its
homotopy exact sequence has zero higher groups, including the
transgression from $\pi_2(S)$. By the cut-annulus calculation in A.6,
the homomorphism on fundamental groups from the link to $T$ is
\[
 \bigl(\Stab(\sigma)\cap\Stab(\tau)\bigr)/\Delta_\tau
       \lhook\joinrel\longrightarrow\Stab(\tau)/\Delta_\tau.
\]
It is injective. Equivalently its restriction to the kernel of
projection to $S$ is the injection of the relative meridian lattice
$\Delta_\sigma/\Delta_\tau$. This injection follows directly from
the independent multitwist coordinates of A.5.

The homotopy exact sequence of the evaluation map now shows that
every component of its homotopy fibre has zero positive homotopy
groups. These spaces have CW homotopy type, so every component is
contractible. Identity-stratum paths give the same conclusion from
asphericity. Hence the exit infinity-category is equivalent to an
ordinary category. Jansen's Corollary 6.6 further identifies this
ordinary category, including its composition, with A.4; Corollary
6.9 gives its stack classifying space. Thus the local criterion and
the explicit global identification have both been used through their
specified comparison maps.
\end{proof}

\begin{result}{A.15}{Proved; unrestricted square refuted}{Exactly when graph normalization and contraction have a square}
Let $j:G'\to G$ be a normalization morphism in the sense of cote 147,
and let $q:G\to G/C$ be a contraction. Put
$C'=j^{-1}(C)\cap E(G')$ and $H'=G'/C'$. A commutative square
\[
\begin{tikzcd}
G'\arrow[r,"j"]\arrow[d]&G\arrow[d,"q"]\\
H'\arrow[r,dashed,"\bar j"]&G/C
\end{tikzcd}
\]
with horizontal normalization morphisms exists exactly when the
following finite tests hold: the induced vertex map is injective,
the genus at every source vertex equals that at its image, and no
retained marked half-edge or leg of $H'$ has been erased by $q$.
The induced maps of the remaining labels must have the incidences
specified in the definition of a normalization morphism; this is
checked directly on the surviving half-edges. When these tests hold,
$\bar j$ is unique, and the associated normalization and relative
generization groupoids form a commuting square up to a canonical
isomorphism. The comparisons are coherent under further admissible
compositions. An unrestricted such square for every pair is false,
already for the strict normalization and smoothing of the single
loop in type $(1,1)$.
\end{result}
\begin{proof}
A commutative square erases from $G'$ precisely the edges whose
images lie in $C$; it therefore forces $C'$. Vertices after contraction
are the connected components of the contracted subgraph. Their
weights are the sum of the original vertex weights plus its first
Betti number. The original vertex injection consequently determines
the candidate vertex map after contraction and its weights. A
normalization morphism requires this map to be injective and to
preserve weights. A marked section retained in $H'$ requires a
surviving target leg or half-edge. These observations prove the
necessity of all the stated tests. Conversely, when the tests hold,
the induced vertex, edge, and label maps satisfy the defining
injection and incidence conditions, so they define $\bar j$.
Surjectivity of the quotient maps on surviving graph data proves
uniqueness.

For the geometric compatibility, use the component description of
A.9 and the equations $x_ey_e=t_e$. The retained nodes have the
same equations in both orders, with parameters indexed by $C'$.
Unselected connected components are discarded. In a contracted
component meeting the selected vertices, injectivity of the new
vertex map says that the selected part is connected after its
own contractions. Equality of the two new genus weights says that
the omitted part consists of genus-zero trees attached to it;
any additional cycle or positive vertex genus would increase the
target weight. The condition on retained sections removes the
sections that would obstruct contraction of those trees. Forgetting
their remaining unselected markings therefore contracts them by
stable reduction. This is exactly the forget-and-stabilize map of
C.1, applied on the selected component. It proves commutation on
families and on the punctured normal germs. Taking fundamental
groupoids gives the asserted square and its endpoint maps.

The comparison restricts to the identity on every retained
normalization component and section. For three operations these
restrictions agree, and uniqueness of stable contraction extends
that agreement to the glued families. This proves coherence.
Disjoint normalization and smoothing of A.10 is the case in which
the tests hold without an omitted attached tree.

For the counterexample take the rational one-loop graph with one
marked leg. Its strict normalization is the single rational vertex
with three legs. Smoothing the loop gives the genus-one vertex with
one leg. The source of a putative lower normalization map still
has genus zero and its only vertex would have to map to a vertex
of genus one. Thus it fails the explicit weight test. This proves
the obstruction by the required comparison data, while the finite
tests above decide every pair for the stated classes of graph maps.
\end{proof}
